\section{Kernel step 2 Restore one word and count its change}
Fill \texttt{video\_scalar\_word}. Load both current input words before storing, compute their XOR, and write the output word. Add one to \texttt{\$v1} only when that XOR word is nonzero. Advance all three pointers by four and decrement words remaining by one. The supplied branch repeats the scalar loop.

The worked first word is unchanged, so its output is zero and the change count is still zero at the first checkpoint. The second word produces hexadecimal \texttt{88} and contributes one change. A scalar complete call produces all five expected words and returns a change count of three.

Loading both words before the store is necessary when output exactly aliases an input. A signed word load can also be correct for this full-word XOR and zero test, provided no high-bit word is mistaken for zero. Do not count only positive signed differences.

\paragraph{Observe} Use \texttt{s02-1}; stop before \texttt{video\_scalar\_stored}. Inspect the state described above before continuing.

\Needspace{7\baselineskip}
\paragraph{Check} What is the change predicate?
\begin{itemize}
\item[A.] The current pixel is nonblack.
\item[B.] The complete XOR word is nonzero.
\item[C.] The XOR word is positive as a signed number.
\end{itemize}
